% page 117 of the facsimile (image p-116 of the scan)
% block 165.1 x 250.6 mm ; left margin 36.9 mm
\pgc{15.240mm}{0.000mm}{
\l{\cel{50}lO7.}
\l{}
\l{\sou{dialogar}. (\sou{netrans., kun}).Konversar kun un persono. -DEFIRS.}
\l{}
\l{\sou{diamanto}. (\sou{mineral.}) Lapido ek karbo kristaligita, un sorto}
\l{\cel{3}di qua uzesas kom juvelo pos facetizo (se lu esas pura).}
\l{\cel{3}DEFIRS.}
\l{}
\l{\sou{diametro}. (\sou{geom.}) Lineo rekta qua pasas tra la centro di cirklo,}
\l{\cel{3}di sfero, di sferoido, e di qua la extremaji tushas la}
\l{\cel{3}periferio.- DEFIRS.}
\l{}
\l{dianto. (\sou{bot.}) Planto gardenala, de la familio \textquotedbl{}kariofilei\textquotedbl{}.}
\l{\cel{3}L.\cel{1}\sou{dianthus}. - L.}
\l{}
\l{\sou{diapazono}. (\sou{muz.}) Instrumento qua donas noto uzata kom departo-}
\l{\cel{3}punto por regular la tono. - DEFIRS.}
\l{}
\l{\sou{diapedezo}. (\sou{patol.})\cel{2}Hemoragio tra la pelo-pori.- DEFIRS.}
\l{}
\l{\sou{diarear}. (\sou{netrans.}) (\sou{patol.}) Ofte evakuar feko liquida ek la}
\l{\cel{3}intestino. - DEFIS.}
\l{}
\l{\sou{diaso}. (\sou{geol.}) Nomo donita olim, da Macon (l839) a la strato}
\l{\cel{3}geologiala permiana. - DE.}
\l{}
\l{\sou{diastazo}. (\sou{kemio)}\cel{1}. Substanco azotoza, neutra, extraktita de}
\l{\cel{3}hordeo, de frumento, e c., qua efikas quale fermento a la ami-}
\l{\cel{3}lo di la grani, tale ke lu separas la sukro e la dextrino de}
\l{\cel{3}la substanci nesolvebla kun qui li esas mixita. - DEFIRS.}
\l{}
\l{\sou{diastolo.(anat.}) Movo di dilato di la kordio e di la arterii,}
\l{\cel{3}qua alternas kun la movo di konstrikto (sistolo). - DEFIS.}
\l{}
\l{\sou{diatezo}. (fiziol.) Ta stando generala di la organismo di homo,}
\l{\cel{3}maxim-multa-kaze atavismala od inata, qua, pos esir latenta,}
\l{\cel{3}su manifestas per ula morbi. - DEFIS.}
\l{}
\l{\sou{diatomeo}. (\sou{bot.}) Familio de algi feoficea, od algi bruna. DEFIRS.}
\l{}
\l{\sou{diatonika.(muziko)}. Qua procedas segun la sucedo naturala di la}
\l{\cel{3}toni e la mi-toni di la gamo. - DEFIRS.}
\l{}
\l{\sou{diatribo}. I.Diserturo kritikoza. - II. Kritiko akra. - DEFIRS.}
\l{}
\l{\sou{dicar}. (\sou{trans., ad, pri, pro}).Savigar per vorti e signi (to}
\l{\cel{3}quon onu havas komunikenda). - FIS.}
\l{}
\l{\sou{dicernar}\cel{1}(\sou{trans.})\cel{2}Vidar o komprenar la difero inter du kozi}
\l{\cel{3}(t. e. ke li interdiferas). - EFIS.}
\l{}
\l{diciono. Arto, maniero parolar o deklamar. - DEFIRS.}
\l{}
\l{dicionario. Vorto-libro en qua esas definita la senco di la}
\l{\cel{3}vorti klasifikita alfabetale (quale en ica radikaro).DEFIRS.}
\l{}
\l{dioiplinar. (\sou{trans.}) Kustumigar a la regulo di aubordino e di}
\l{\cel{3}drdiRb, imobiata a la membri di korpb, di organismo. - DEFIES.}
}
